\documentclass[10pt,a4paper]{amsart}
\usepackage[margin=19mm]{geometry}
\usepackage{amsmath,amssymb,mathtools}
\usepackage[scr=rsfs,cal=euler]{mathalfa}
\usepackage{xfrac}
\usepackage[unicode,hidelinks,bookmarks=false]{hyperref}
\newcommand{\ie}{i.e.\ }
\newcommand{\cf}{cf.\ }
\newcommand{\resp}{respectively\ }
\newcommand{\Subsection}{\subsection}
\providecommand{\nobreakdash}{\nobreak\hbox{-}}
\setlength{\emergencystretch}{2em}
\multlinegap0pt
\usepackage{bm}
\usepackage{bbm}
\usepackage{dsfont}
\usepackage{xspace}
\usepackage{cancel}
\usepackage{mathtools}
\usepackage{amsthm}
\makeatletter
\@ifundefined{Theorem}{\newtheorem{Theorem}{Theorem}}{}
\@ifundefined{Lemma}{\newtheorem{Lemma}[Theorem]{Lemma}}{}
\makeatother
\newcommand{\Om}{\Omega}
\newcommand{\mR}{\mathbb{R}}                    % Formatting for R
\newcommand{\norm}[1]{\lVert #1 \rVert}         % Formatting for the norm
\newcommand{\ol}[1]{\overline{#1}}
\newcommand{\p}{\partial}
\title[Inverse problems for semilinear elliptic PDE with a general ]{Inverse problems for semilinear elliptic PDE with a general nonlinearity $a(x,u)$}
\author{David Johansson, Janne Nurminen, Mikko Salo}
\date{Source: arXiv:2312.12196v2 (CC BY 4.0)}
\begin{document}
\maketitle

\noindent\emph{Source and licence.} Inverse problems for semilinear elliptic PDE with a general nonlinearity $a(x,u)$. Version arXiv:2312.12196v2 (CC BY 4.0), distributed under the Creative Commons Attribution 4.0 International licence, as recorded in the arXiv licence metadata for this exact version and shown on the abstract page https://arxiv.org/abs/2312.12196v2. Full rights notice: licenses/arxiv-2312.12196-CC-BY-4.0.txt.\par\smallskip

\noindent\textit{Editorial scope.} Counted unit: the Lemma whose source label is \texttt{lemma:inductive\_higher\_linearization} in the article (source0.tex lines 1219--1241, source label \texttt{lemma:inductive\_higher\_linearization}), delivered together with the article's own complete original proof of it (source0.tex lines 1242--1319, one \texttt{proof} environment of 78 lines). No mathematics is written, repaired or re-derived: statement and proof are the article's own text line for line. Required context retained uncounted: lemma\_perturbed\_sol (lines 610--635); lemma:smooth\_solution\_operator (lines 1013--1019). Every definition, display and label of those lines is retained unchanged. Omitted from this package and not redistributed: 1--609, 636--1012, 1020--1218, 1320--1518. Nothing inside a retained excerpt is omitted or altered: every retained body is delivered exactly as the article's lines are written, so no omission receipt of any kind accompanies this package. Citations: the retained excerpts carry no citation command at all, so no bibliography entry of the article is transcribed and no reference is redistributed. Reference adaptation: none was needed and none was made; every reference inside the delivered unit resolves inside it, so no unresolved reference or citation is produced. Printed slips kept literal: no printed word, symbol, hypothesis or number of the counted statement or of its proof was altered. Layout and class: the article's own class is \texttt{amsart} with packages including fontenc, mathtools, amsmath, amscd, amssymb, amsthm, comment, enumitem, graphicx, mathrsfs, hyperref, biblatex, calc, xcolor; the wrapper is the lane's pinned \texttt{amsart} editorial wrapper at 10pt on A4, which restates the theorem-like environments the retained text needs and adds the article's own macro definitions that the retained text needs (\texttt{\textbackslash Om}, \texttt{\textbackslash mR}, \texttt{\textbackslash norm}, \texttt{\textbackslash ol}, \texttt{\textbackslash p}), with extra wrapper packages (bm, bbm, dsfont, xspace, cancel, mathtools). The delivered numbering is the unit's own. Assets: no retained span contains a figure environment, a graphics-inclusion command, a PDF-page inclusion, a table or a TikZ picture; no image file is copied into this package, the package declares no asset dependency and no image file is redistributed with it. Licence: version arXiv:2312.12196v2 of this article is distributed under the Creative Commons Attribution 4.0 International licence, as recorded in the arXiv licence metadata for this exact version and shown on the abstract page https://arxiv.org/abs/2312.12196v2. A full rights notice is delivered at licenses/arxiv-2312.12196-CC-BY-4.0.txt. Characters: every retained excerpt is pure ASCII, so the delivered engine, XeLaTeX with its default Latin Modern text font, reports no missing character.
\begin{Lemma} \label{lemma_perturbed_sol}
Let $a \in C^{k}(\mR, C^{1,\alpha}(\ol{\Om}))$, $k \geq 3$.
Let $w \in C^{2,\alpha}(\ol{\Omega})$ be a solution of $\Delta w + a(x,w) = 0$.
Let $q(x) = \p_u a(x,w(x))$.
Then there exist $\delta, C > 0$ and a $C^{k-1}$ map $Q = Q_{a,w}: B_{\delta} \to B_{\delta}$ satisfying 
\begin{align*}
Q(B_{\delta}) &\subseteq N_q^{\perp}, \\
Q(B_{\delta})|_{\p \Omega} &\subseteq \p_{\nu} N_q, \\
Q(0) &= DQ(0) = 0
\end{align*}
and
\begin{equation}\label{lemma_estimate_r}
    \norm{Q(v)}_{C^{2,\alpha}(\ol{\Omega})}\leq C\norm{v}_{C^{2,\alpha}(\ol{\Omega})}^2,
\end{equation}
such that \( S_{a}\colon B_{\delta}\to C^{2,\alpha}(\ol{\Omega}) \) defined by \( u = S_{a,w}(v) = w + v + Q(v) \) is a \( C^{k-1} \) map satisfying 
\begin{equation} \label{u_v_eq}
	\Delta u + a(x,u) = \Delta v + q v \quad\text{in }\Omega
\end{equation}
with \( S_{a,w}'(0)v = v \). In particular, if $v$ solves $\Delta v + qv = 0$, then $u = S_{a,w}(v)$ solves $\Delta u + a(x,u) = 0$.

Conversely, if $\delta$ is small enough, then given any solution $u \in C^{2,\alpha}(\ol{\Omega})$ of $\Delta u + a(x,u) = 0$ with $\norm{u-w}_{C^{2,\alpha}(\ol{\Omega})} \leq \delta$ there exists a unique solution $v \in C^{2,\alpha}(\ol{\Omega})$ of $\Delta v + qv = 0$ such that $u = S_{a,w}(v)$. The function $v$ is explicitly given by 
\begin{equation} \label{u_v_converse}
v = P_{N_q}(u-w) + G_q(0, (u-w)|_{\p \Omega}),
\end{equation}
and one has $\norm{v}_{C^{2,\alpha}(\ol{\Omega})} \leq C \norm{u-w}_{C^{2,\alpha}(\ol{\Omega})}$.
\end{Lemma}

\begin{Lemma}\label{lemma:smooth_solution_operator}
Let $a_1,a_2\in C^{k+1}(\mR, C^{2,\alpha}(\ol{\Om}))$ with $k\geq 2$ and let $w_1,w_2$ have the same Cauchy data and solve $\Delta w_i + a_i(x,w_i) = 0$ in $\Omega$.
Write $q_i = \p_u a_i(x,w_i)$. Let $S_{a_1}\colon V_{q_1,\delta_1}\to C^{2,\alpha}(\ol{\Om})$ be the solution map from Lemma \ref{lemma_perturbed_sol}, for some $\delta_1>0$.
Suppose $u_{1,v} = S_{a_1}(v)$ and that $C_{a_1}^{w_1,\delta} \subseteq C_{a_2}^{0,C}$.
Then there exists a $\delta_2 > 0$ and a $C^k$ map $T_{a_2}\colon V_{q_1,\delta_2}\to C^{2,\alpha}(\ol{\Om})$, $T_{a_2}(v) = u_{2,v}$, where $u_{2,v}$ has the same Cauchy data as $u_{1,v}$ and solves $\Delta u_{2,v} + a_2(x,u_{2,v}) = 0$.
Moreover, when \( \partial_{u}a_{1}(x,w_{1}) = \partial_{u}a_{2}(x,w_{2}) \) then $T_{a_2}'(0)v = v$.
\end{Lemma}

\begin{Lemma}\label{lemma:inductive_higher_linearization}
Let \( a_{1},a_{2} \in C^{k+2}(\mathbb{R},C^{2,\alpha}(\ol{\Omega}))\) with $k\geq 1$.
Let \( S_{a_{1}},T_{a_2} \) be the solution operators \( \Delta u + a_{i}(x,u)=0 \) from Lemma \ref{lemma_perturbed_sol} and Lemma \ref{lemma:smooth_solution_operator}.
Denote by \( R(v)\coloneqq S_{a_{1}}(v)-T_{a_{2}}(v) \) and let \( u = D^{l}R(0; v_{1},\ldots, v_{k}) \) for \( 1 \leq l \leq k \).
Then \( u \) satisfies
\begin{equation}\label{eq-inductive_higher_linearization_new_1}
\begin{cases}
	[\Delta + \partial_{u}a_{2}(x,w_{2})]u = F_{l} + [\partial_{u}^{l}a_{2}(x,w_{2})-\partial_{u}^{l}a_{1}(x,w_{1})]\prod_{i=1}^{l}h_{i}&\text{in }\Omega, \\
	u = \partial_{\nu}u = 0&\text{on }\partial\Omega.
\end{cases}
\end{equation}
where \( h_{i} = S_{a_{1}}^{\prime}(0;v_{i}) \), \( w_{1} = S_{a_{1}}(0) \), \( w_{2} = T_{a_{2}}(0) \), and
\begin{equation*}
	F_{l} = \sum_{i=2}^{l}\partial_{u}^{i}a_{2}(x,w_{2})P_{i}+\sum_{i=1}^{l-	1}[\partial_{u}^{i}a_{2}(x,w_{2})-\partial_{u}^{i}a_{1}(x,w_{1})]Q_{i}
\end{equation*}
whenever $2\leq l\leq k$, and otherwise $F_1 = 0$.
Here \( P_{i} \) is a polynomial in \( D^{j}S_{a_{1}}(0;v_{1},\ldots, v_{j}) \), \( D^{j}T_{a_{1}}(0;v_{1},\ldots, v_{j}) \), and \( D^{j}R(0;v_{1},\ldots, v_{j}) \) and \( Q_{i} \) is a polynomial in \( D^{j}S_{a_{1}}(0;v_{1},\ldots, v_{j}) \) for \( 1\leq j\leq i \).
Both \( P_{i} \) and \( Q_{i} \) vanish at the origin and each term in \( P_{i} \) contains at least one factor which is \( D^{j}R(0;v_{1},\ldots,v_{j}) \).
Additionally, \( F_{l} \) satisfies
\begin{equation}\label{eq-inductive_higher_linearization_new_2}
	\partial_{u}^{i}a_{1}(x,w_{1}) = \partial_{u}^{i}a_{2}(x,w_{2})\text{ for }1\leq i\leq l-1\implies F_{l} = 0.
\end{equation}
\end{Lemma}

\begin{proof}
Since \( S_{a_{1}}(v),T_{a_{2}}(v) \) have equal Cauchy data, \( R(v) \) has vanishing Cauchy data for every \( v \).
Then the derivatives of \( R(v) \) also have vanishing Cauchy data, which establishes the boundary condition in \eqref{eq-inductive_higher_linearization_new_1}.

Deriving the linearized differential equation is somewhat more involved.
The first linearization is
\begin{equation}\label{eq-inductive_higher_linearization_new_3}
	[\Delta+\partial_{u}a_{2}(x,w_{2})]u = [\partial_{u}a_{2}(x,w_{2})-\partial_{u}a_{1}(x,w_{1})]h_{1}.
\end{equation}
In other words, \( F_{1} = 0 \).
The second linearization is
\begin{equation*}
\begin{aligned}
	[\Delta+\partial_{u}a_{2}(x,w_{2})]u = &-\partial_{u}^{2}a_{2}(x,w_{2})[S_{a_{1}}^{\prime}(0;v_{1})R^{\prime}(0;v_{2}) + T_{a_{1}}^{\prime}(0;v_{2})R^{\prime}(0;v_{1})] \\
	&+[\partial_{u}a_{2}(x,w_{2})-\partial_{u}a_{1}(x,w_{1})]S_{a_{1}}^{\prime\prime}(0;v_{1},v_{2}) \\
	&+ [\partial_{u}^{2}a_{2}(x,w_{2})-\partial_{u}^{2}a_{1}(x,w_{1})]h_{1}h_{2}.
\end{aligned}
\end{equation*}
Recall that \( h_{i} = S_{a_{1}}^{\prime}(0,v_{i}) \).
We see that
\begin{equation*}
\begin{aligned}
	F_{2} = &-\partial_{u}^{2}a_{2}(x,w_{2})[S_{a_{1}}^{\prime}(0;v_{1})R^{\prime}(0;v_{2}) + T_{a_{1}}^{\prime}(0;v_{2})R^{\prime}(0;v_{1})] \\
	&+[\partial_{u}a_{2}(x,w_{2})-\partial_{u}a_{1}(x,w_{1})]S_{a_{1}}^{\prime\prime}(0;v_{1},v_{2}).
\end{aligned}
\end{equation*}
We proceed by induction.
Suppose that \( \tilde{u} = D^{l-1}R(0;v_{1},\ldots,v_{l-1}) \) satisfies
\begin{equation*}
	[\Delta + \partial_{u}a_{2}(x,w_{2})]\tilde{u} = F_{l-1} + [\partial_{u}^{l-1}a_{2}(x,w_{2})-\partial_{u}^{l-1}a_{1}(x,w_{1})]\prod_{i=1}^{l-1}h_{i},
\end{equation*}
and denote again by \( u \) the derivative \( D^{l}R(0,v_{1},\ldots,v_{l}) \).
We linearize this term by term.
The derivative of \( (\Delta+\partial_{u}a_{2}(x,w_{2}))\tilde{u} \) at \( 0 \) and in direction \( v_{l} \) is
\begin{equation*}
	(\Delta+\partial_{u}a_{2}(x,w_{2}))u + \partial_{u}^{2}a_{2}(x,w_{2})T_{a_{2}}^{\prime}(0;v_{l})\tilde{u},
\end{equation*}
and the latter term will be incorporated into \( F_{l} \).
In the following computations \( D_{v} \) denotes the derivative at \( 0 \) in direction \( v \).
The derivative of \( F_{l-1} \) at \( 0 \) and in direction \( v_{l} \) is
\begin{equation*}
\begin{aligned}
	&\sum_{i=2}^{l-1}[\partial_{u}^{i+1}a_{2}(x,w_{2})T_{a_{2}}^{\prime}(0;v_{l})P_{i} + \partial_{u}^{i}a_{2}(x,w_{2})D_{v_{l}}P_{i}] \\
	&+\sum_{i=1}^{l-2}[\partial_{u}^{i+1}a_{2}(x,w_{2})\underbrace{T_{a_{2}}^{\prime}(0;v_{l})}_{=S_{a_{1}}^{\prime}(0;v_{l})-R^{\prime}(0;v_{l})}-\partial_{u}^{i+1}a_{1}(x,w_{1})S_{a_{1}}^{\prime}(0;v_{l})]Q_{i} \\
	&+\sum_{i=1}^{l-2}[\partial_{u}^{i}a_{2}(x,w_{2})-\partial_{u}^{i}a_{1}(x,w_{1})]D_{v_{l}}Q_{i} \\
	&=\sum_{i=2}^{l}\partial_{u}^{i}a_{2}(x,w_{2})[(1-\delta_{i,2})(T_{a_{2}}^{\prime}(0;v_{l})P_{i-1} - R^{\prime}(0;v_{l})Q_{i-1}) + (1-\delta_{i,l})D_{v_{l}}P_{i}] \\
	&+\sum_{i=1}^{l-1}[\partial_{u}^{i}a_{2}(x,w_{2})-\partial_{u}^{i}a_{1}(x,w_{1})][(1-\delta_{i,1})S_{a_{1}}^{\prime}(0;v_{l})Q_{i-1} + (1-\delta_{i,l-1})D_{v_{l}}Q_{i}],
\end{aligned}
\end{equation*}
where \( \delta_{i,j} \) denotes the Kronecker delta.
Lastly, we differentiate \( [\partial_{u}^{l-1}a_{2}(x,w_{2})-\partial_{u}^{l-1}a_{1}(x,w_{1})]\prod_{i=1}^{l-1}h_{i} \) at 0 and in direction \( v_{l} \).
We get
\begin{equation*}
\begin{aligned}
	&[\partial_{u}^{l}a_{2}(x,w_{2})\underbrace{T_{a_{2}}^{\prime}(0;v_{l})}_{=S_{a_{1}}^{\prime}(0;v_{l})-R^{\prime}(0;v_{l})}-\partial_{u}^{l}a_{1}(x,w_{1})S_{a_{1}}^{\prime}(0;v_{l})]\prod_{i=1}^{l-1}h_{i} \\
	&+ [\partial_{u}^{l-1}a_{2}(x,w_{2})-\partial_{u}^{l-1}a_{1}(x,w_{1})]D_{v_{l}}\Big(\prod_{i=1}^{l-1}h_{i}\Big) \\
	&=[\partial_{u}^{l-1}a_{2}(x,w_{2})-\partial_{u}^{l-1}a_{1}(x,w_{1})]D_{v_{l}}\Big(\prod_{i=1}^{l-1}h_{i}\Big) -\partial_{u}^{l}a_{2}(x,w_{2})R^{\prime}(0;v_{l})\prod_{i=1}^{l-1}h_{i}\\
	& +[\partial_{u}^{l}a_{2}(x,w_{2})-\partial_{u}^{l}a_{1}(x,w_{1})]\prod_{i=1}^{l}h_{i}.
\end{aligned}
\end{equation*}
Here, the terms on the second to last row will be incorporated into \( F_{l} \) and the last row is not incorporated into \( F_{l} \).
Now let
\begin{equation*}
	\tilde{P}_{i} = (1-\delta_{i,2})(T_{a_{2}}^{\prime}(0;v_{l})P_{i-1} - R^{\prime}(0;v_{l})Q_{i-1}) + (1-\delta_{i,l})D_{v_{l}}P_{i} - \delta_{i,2}T_{a_{2}}^{\prime}(0;v_{l})\tilde{u}-\delta_{il}R^{\prime}(0;v_{l})\prod_{i=1}^{l-1}h_{i}
\end{equation*}
and
\begin{equation*}
\begin{aligned}
	\tilde{Q}_{i} = [(1-\delta_{i,1})S_{a_{1}}^{\prime}(0;v_{l})Q_{i-1} + (1-\delta_{i,l-1})D_{v_{l}}Q_{i}] + (1-\delta_{i,l-1})D_{v_{l}}\Big(\prod_{i=1}^{l-1}h_{i}\Big).
\end{aligned}
\end{equation*}
By using \( \tilde{P}_{i} \), \( \tilde{Q}_{i} \) in place of \( P_{i} \), \( Q_{i} \) in the definition of \( F_{l} \), we find that the linearization of order \( l \) indeed is \eqref{eq-inductive_higher_linearization_new_1}.

A careful investigation of the above procedure reveals that every term in \( P_{i} \) indeed must contain at least one derivative \( D^{j}R \).
To see that \eqref{eq-inductive_higher_linearization_new_2} holds, first note that if \( \partial_{u}a_{1}(x,w_{1}) = \partial_{u}a_{2}(x,w_{2}) \) then the right-hand side of the first linearization \eqref{eq-inductive_higher_linearization_new_3} vanishes and it follows by unique continuation that \( R^{\prime}(0;v_{j}) = 0 \), since it has vanishing Cauchy data.
Then every term in \( F_{2} \) has a factor which vanishes, hence \( F_{2} = 0 \).
We obtain \eqref{eq-inductive_higher_linearization_new_2} by iterating this procedure to higher order linearizations.
\end{proof}

\end{document}
